\section{Divisibility of the minor}\label{sec:arithmetic}
We write $q$ for a general prime, reserving $p$ for the auxiliary prime in
Section~\ref{sec:rank}. The basic Frobenius argument in Section~\ref{subsec:model} gives a row
contribution from each inert prime and a loss depending on the derivative
orders selected. We now strengthen the row contribution, first by using the
conjugate variable and then by retaining the small primes. Finally we explain
how class number one controls the losses from non-inert primes.

Recall that for an inert prime $q\geq K$, the integral row basis is
$P_{r,s}=z^{r\bmod q}(z^q-w)^{\lfloor r/q\rfloor}w^s$.
Writing $v_q$ for the valuation normalized by $v_q(q)=1$, the calculation
in Lemma~\ref{lem:frobenius} gives
\begin{equation}\label{eq:basic-valuation}
 v_q(\partial_z^{[j]}P_{r,s}(z_i,\bar z_i))
 \geq\max(0,\lfloor r/q\rfloor-\lfloor j/q\rfloor).
\end{equation}
Each determinant term uses every row and every selected column once. Hence
\begin{equation}\label{eq:basic-minor-valuation}
 v_q(\Delta_C)\geq\sum_\Lambda\lfloor r/q\rfloor
                  -\sum_C\lfloor j/q\rfloor.
\end{equation}
The right side can be negative, in which case it is merely a weak bound.
Its usefulness emerges after summing over primes: the column loss combines
with the factorials in the Hasse derivatives, as in \eqref{eq:F-H}.

\subsection{Additional divisibility in the conjugate variable}
Differentiation in $z$ leaves polynomials in $w$ untouched. Such polynomials can contribute additional divisibility at no derivative cost. We use the following elementary tower. Define
\[
 T_1(w)=w^{q^2}-w,\qquad
 \nu_a=\frac{q^{2a}-1}{q^2-1},\qquad
 T_{a+1}(w)=T_a(w)^{q^2}-q^{\nu_a(q^2-1)}T_a(w).
\]
These are monic polynomials of degrees $q^{2a}$. Inductively,
\begin{equation}\label{eq:tower-value}
 T_a(\bar z_i)\in q^{\nu_a}\OO_D.
\end{equation}
Indeed, if $T_a(\bar z_i)=q^{\nu_a}b$, then the next value equals
$q^{q^2\nu_a}(b^{q^2}-b)$, which is divisible by
$q^{q^2\nu_a+1}=q^{\nu_{a+1}}$.

Write $s=\sum_{a\geq0}d_aq^{2a}$ in base $q^2$, and put
\[
 A_s(w)=w^{d_0}\prod_{a\geq1}T_a(w)^{d_a}.
\]
It is monic of degree $s$, and \eqref{eq:tower-value} gives
\[
 v_q(A_s(\bar z_i))\geq
 \sum_{a\geq1}d_a\nu_a
 =\sum_{a\geq1}\floor{s/q^{2a}}.
\]
Replacing $w^s$ by $A_s(w)$ in the previous basis preserves triangularity. We obtain the stronger row estimate
\begin{equation}\label{eq:tower-valuation}
 v_q\bigl(\partial_z^{[j]}(z^{r\bmod q}(z^q-w)^{\floor{r/q}}A_s(w))
                 (z_i,\bar z_i)\bigr)
 \geq\max(0,\floor{r/q}-\floor{j/q})
         +\sum_{a\geq1}\floor{s/q^{2a}}.
\end{equation}
This is a local statement for any nodes in the order; it has no rank hypothesis.

The row contributions are finite sums over the triangle. Their exact
floor-sum formulas are given in the supplementary notes, Section~\ref{supp:rows}.

\subsection{Small primes and non-inert primes}
For an inert prime $q<K$, replacing $z^q$ by $w$ increases weighted degree. The number of replacements in a row is therefore bounded by
\begin{equation}\label{eq:small-h}
 h_q(r,s)=\min\left(\floor{r/q},
          \floor{(R-1-r-Ks)/(K-q)}\right).
\end{equation}
There is still a substantial contribution: summing this function over the triangle yields
\[
 2\sum_{\Lambda}h_q(r,s)=W_2=\frac{R(S^2-1)}3.
\]
One way to verify the identity is to sum the sizes of its superlevel sets. The condition $h_q\geq a$ is equivalent, after writing $r=qa+r'$, to
$r'+Ks<R-Ka$, and its number of solutions is $K(S-a)(S-a+1)/2$. Summing for $1\leq a<S$ gives the displayed value.

The small-prime estimate used in the computation improves this elementary bound by combining $z^q-w$, $w^q-z$, and $z^{q^2}-z$. The first factor can lose a power of $q$ only after spending at least $q$ orders of differentiation; the latter two can lose one after spending one order. Optimizing where each type is used is a finite triangular row calculation. The resulting finite lower bound is recorded in the supplementary notes, Section~\ref{app:arithmetic}, together with its cut parameter. In particular, $q=2$ is included. Working in the maximal order is essential here, since $2$ is inert when $D\equiv3\pmod8$.

The arithmetic calculation also uses primes above $(D-3)/4$, where inertness is no longer guaranteed. Class number one still supplies useful information. If $q$ splits or ramifies, an ideal of norm $q$ is principal, so
\begin{equation}\label{eq:prime-norm}
 q=\N(x_0+y\omega)=\left(x_0+\frac y2\right)^2+\frac D4y^2,
 \qquad y\ne0.
\end{equation}
At a fixed norm cutoff there are only finitely many such pairs. One can therefore start with an all-prime estimate and subtract an upper bound for the contribution of possible non-inert primes, using these norm representations. This deduction uses principality; it does not test the quadratic character at a particular discriminant.

Counting all norm values as possible primes overestimates the loss. Unique factorization allows some of those values to be removed in advance. Products of two or three nonrational factors have composite norm, and the corresponding primitive products can be counted with their multiplicities controlled. The product counts used for this subtraction are recorded in the supplementary arithmetic notes. At the finite comparisons in Section~\ref{sec:completion}, only the two- and three-factor corrections are used. The upper bound for possible losses is decreased by a rigorously bounded contribution of these composite values. Thus this correction \emph{increases} $\arith$; it is not an additional subtraction of prime divisibility.

For orientation, the resulting arithmetic term has the form
\begin{equation}\label{eq:arithmetic-parts}
 \arith=A_{\rm all}-L_{\rm norm}+C_2+C_3.
\end{equation}
The first term sums inert-prime row credits as a provisional all-prime bound; $L_{\rm norm}$ pays for norm values that might be non-inert primes; $C_2$ and $C_3$ restore credits paid unnecessarily for certified composite values. Split and ramified primes also have their own local conjugate-variable contributions. These, and the derivative-order refunds, are included in the arithmetic lower bound. This matters numerically near $D=499$.

\subsection{The arithmetic estimate used below}
We use the following consequence of the local estimates, in a form that
continues to hold when the number of nodes is reduced. The normalization is $\log|\Delta_C|^2/W$: all prime-sum and
norm-loss terms in $\arith$ are divided by $W$, while $b(j)$ is the column
debit in logarithmic units before this division.

\begin{proposition}[Local arithmetic estimate]\label{prop:arithmetic}
Let $-D$ be a fundamental discriminant with $h(-D)=1$ at one of the
finite comparisons below $875$ in Section~\ref{sec:completion}, with the
parameters specified there. Let $z_1,\dots,z_n$ be the $n$ nearest elements
of $\OO_D$, and let $V$ be the matrix~\eqref{eq:jet-matrix} with
$0\leq j<J$ and $\theta=M/(nJ)\leq1$. Define $\arith,\sigma,c,b(j)$ by
the local arithmetic formulas in the supplementary notes,
Section~\ref{app:arithmetic}, including the norm-product corrections
and derivative-order refunds.
Then~\eqref{eq:arithmetic-interface} and~\eqref{eq:column-remainder} hold for
every set $C$ of $M$ columns with $\Delta_C\ne0$.
\end{proposition}

Here $\Delta_C$ is a minor of the original integral Hasse matrix $V$.
The bound on $D$ will follow by comparing these local estimates with the size
estimate of Section~\ref{sec:metric}. The preceding arguments explain the
constructions behind the estimates. The supplementary notes give the finite formulas and their derivations. Every quantity
depending on $n$ is evaluated at the node count in use. This includes the
number of columns of each order, and hence the least possible sum of their
first-power floors and the largest possible remainder sum.

The higher-prime-power saving can be seen directly from Legendre's formula:
\[
 \log(j!)=\sum_q\sum_{a\geq1}\floor{j/q^a}\log q.
\]
The debit $\sum_q\floor{j/q}\log q$ from \eqref{eq:basic-minor-valuation} cancels its $a=1$ part exactly. The remaining terms, together with the split-prime refunds, form $F(j)$ in \eqref{eq:F-H}. Sorting $H(0),\ldots,H(J-1)$ in decreasing order, say as $H_{(1)},\ldots,H_{(J)}$, and writing $M=nt+r$ with $0\leq r<n$, yields the exact capacity bound
\begin{equation}\label{eq:sorted-remainder}
 \sum_C H(j)\leq n\sum_{a=1}^{t}H_{(a)}+rH_{(t+1)}.
\end{equation}
When $r=0$ the final term is omitted. The right side gives an admissible $Mc/2$. Grouping the floor functions into intervals on which their quotients are
constant gives a finite evaluation of the right side. The resulting bound
for $c$ includes the remainder at every derivative order.
